\documentclass[11pt]{article}

\usepackage[T1]{fontenc}
\usepackage[margin=1in]{geometry}
\usepackage{amsmath, amssymb, amsthm}
\usepackage{enumitem}
\usepackage{hyperref}
\usepackage{booktabs}
\usepackage{listings}
\usepackage{xcolor}

\lstset{
  basicstyle=\ttfamily\small,
  frame=single,
  breaklines=true,
  columns=fullflexible,
  mathescape=true
}

\theoremstyle{definition}
\newtheorem{theorem}{Theorem}[section]
\newtheorem{definition}[theorem]{Definition}
\newtheorem{example}[theorem]{Example}

\title{CS 250 --- Intermezzo: More on Natural Induction}
\author{}
\date{}

\begin{document}
\maketitle

To illustrate natural induction, let's work through some proofs similar to those in The Natural Number Game, but informally.

For the purposes of these proofs we'll assume that the definition of addition by recursion on the second argument is
\begin{align*}
n + 0 &= n \\
n + (\text{succ}\; m) &= \text{succ}\;(n + m)
\end{align*}

\section{Proof that 0 + m = m}
We have automatically from the definition of addition by recursion that $n + 0 = n$, but what about the other way? You might wonder why this needs a proof at all---isn't it obvious? The subtlety is that our definition of addition recurses on the \emph{second} argument. The equation $n + 0 = n$ is the base case of that recursion, so it holds by definition. But $0 + m = m$ puts zero in the \emph{first} argument, which is not the recursive position. The definition doesn't directly tell us what happens there, so we need induction.

\textbf{Theorem:} $\forall m.\; 0 + m = m$

\textbf{Proof} by induction on $m$.

\textbf{Case $m = 0$:} Need to prove that $0 + 0 = 0$, but this follows immediately by the definition of addition.

\textbf{Case $m = \text{succ}\; m'$:}

Assuming $0 + m' = m'$.

Shows $0 + \text{succ}\; m' = \text{succ}\; m'$.

By the definition of addition,
\begin{align*}
0 + \text{succ}\; m' &= \text{succ}\;(0 + m')
\end{align*}
but
\begin{align*}
\text{succ}\;(0 + m') &= \text{succ}\; m' && \text{by the induction hypothesis that } 0 + m' = m'
\end{align*}

QED

\section{Proof that succ n + m = succ (n + m)}

\textbf{Theorem:} $\forall n\, m.\; \text{succ}\; n + m = \text{succ}\;(n + m)$

\textbf{Proof} by induction on $m$.

\textbf{Case $m = 0$:} $\text{succ}\; n + 0 = \text{succ}\;(n + 0)$.

Immediate from the definition of addition.

\textbf{Case $m = \text{succ}\; m'$:}

Assuming: $(\text{succ}\; n) + m' = \text{succ}\;(n + m')$.

Shows: $(\text{succ}\; n) + (\text{succ}\; m') = \text{succ}\;(n + (\text{succ}\; m'))$.

\begin{align*}
(\text{succ}\; n) + (\text{succ}\; m')
  &= \text{succ}\;(\text{succ}\; n + m') && \text{by the definition of addition} \\
  &= \text{succ}\;(\text{succ}\;(n + m')) && \text{by the induction hypothesis}
\end{align*}

but the right hand side tells us
\begin{align*}
\text{succ}\;(n + (\text{succ}\; m')) &= \text{succ}\;(\text{succ}\;(n + m')) && \text{by the definition of addition}
\end{align*}

so we end up with
\[
\text{succ}\;(\text{succ}\;(n + m')) = \text{succ}\;(\text{succ}\;(n + m'))
\]
which is true by reflexivity.

QED

\section{Proof that addition is associative}\index{associativity}

\textbf{Theorem:} $\forall n\, m\, r.\; n + (m + r) = (n + m) + r$

\textbf{Proof} by induction on $n$.

\textbf{Case $n = 0$:} Need to prove that $\forall m\, r.\; 0 + (m + r) = (0 + m) + r$.

Both sides reduce to $m + r$ by the left and right identity properties of 0 (proved in Section 1 and given by definition, respectively):
\[
m + r = m + r
\]
which is true by reflexivity.

\textbf{Case $n = \text{succ}\; n'$:} Need to prove that \textbf{assuming} $\forall m\, r.\; n' + (m + r) = (n' + m) + r$ we can derive $\text{succ}\; n' + (m + r) = (\text{succ}\; n' + m) + r$.

\begin{align*}
\text{succ}\; n' + (m + r) &= \text{succ}\;(n' + (m + r)) && \text{by our theorem above}
\end{align*}

and the right hand side similarly transforms
\begin{align*}
(\text{succ}\; n' + m) + r
  &= \text{succ}\;(n' + m) + r && \text{by our theorem above} \\
  &= \text{succ}\;((n' + m) + r) && \text{by the same theorem} \\
  &= \text{succ}\;(n' + (m + r)) && \text{by the induction hypothesis}
\end{align*}

so we have
\[
\text{succ}\;(n' + (m + r)) = \text{succ}\;(n' + (m + r))
\]
which is true by reflexivity.

QED

\section{Proof that addition is commutative}\index{commutativity}

\textbf{Theorem:} $\forall n\, m.\; n + m = m + n$

\textbf{Proof} by induction on $n$.

\textbf{Case $n = 0$:} Need to prove that $\forall m.\; 0 + m = m + 0$.

This is immediately solved by the fact that 0 is the left and right identity for addition, so this simplifies to $m = m$ which is true by the definition of $=$ (i.e.\ by reflexivity).

\textbf{Case $n = \text{succ}\; n'$:} Need to prove that \textbf{assuming} $\forall m.\; n' + m = m + n'$ we can show $\forall m.\; (\text{succ}\; n') + m = m + (\text{succ}\; n')$.

So by the definition of addition we have
\[
m + (\text{succ}\; n') = \text{succ}\;(m + n')
\]

and by one of the above theorems we have that
\[
(\text{succ}\; n') + m = \text{succ}\;(n' + m)
\]
but
\[
\text{succ}\;(n' + m) = \text{succ}\;(m + n') \quad \text{by the induction hypothesis}
\]

and thus
\[
\text{succ}\;(m + n') = \text{succ}\;(m + n')
\]
which is true by reflexivity.

QED

\end{document}
